\chapter{2012年湖南卷（文）}

\section{单选题}

\begin{problem}
设集合 \(M = \{-1, 0, 1\}\)，\(N = \{x \mid x^2 = x\}\)，则 \(M \cap N =\)（\quad）

\choices
    {\( \{-1,0,1\} \)}
    {\(\{0, 1\}\)}
    {\( \{1\} \)}
    {\( \{0\} \)}
\end{problem}

\begin{answer}
B
\end{answer}

\begin{solution}
由 \(x^2=x\) 得 \(x(x-1)=0\)，所以 \(N=\{0,1\}\). 因而
\(M\cap N=\{-1,0,1\}\cap\{0,1\}=\{0,1\}\)，故选 B.
\end{solution}

\begin{problem}
复数 \(z = \i(\i + 1)\)（\(\i\) 为虚数单位）的共轭复数是（\quad）

\choices
    {\(-1 - \i\)}
    {\(-1 + \i\)}
    {\(1 - \i\)}
    {\(1 + \i\)}
\end{problem}

\begin{answer}
A
\end{answer}

\begin{solution}
先将复数化为代数形式：
\(z=\i(\i+1)=\i^2+\i=-1+\i\). 共轭复数是实部不变、虚部变号，故
\(\overline z=-1-\i\)，所以选 A.
\end{solution}

\begin{problem}
命题“若 \(\alpha = \frac{\pi}{4}\)，则 \(\tan \alpha = 1\)”的逆否命题是（\quad）

\choices
    {若 \(\alpha \neq \frac{\pi}{4}\)，则 \(\tan \alpha \neq 1\)}
    {若 \(\alpha = \frac{\pi}{4}\)，则 \(\tan \alpha \neq 1\)}
    {若 \(\tan \alpha \neq 1\)，则 \(\alpha \neq \frac{\pi}{4}\)}
    {若 \(\tan \alpha \neq 1\)，则 \(\alpha = \frac{\pi}{4}\)}
\end{problem}

\begin{answer}
C
\end{answer}

\begin{solution}
设命题中的前件为 \(p\)：\(\alpha=\frac{\pi}{4}\)，后件为 \(q\)：\(\tan\alpha=1\). 命题
“若 \(p\)，则 \(q\)”的逆否命题是“若非 \(q\)，则非 \(p\)”，即
“若 \(\tan\alpha\neq1\)，则 \(\alpha\neq\frac{\pi}{4}\)”. 故选 C.
\end{solution}

\begin{problem}
某几何体的正视图和侧视图均如下图所示，则该几何体的俯视图不可能是（\quad）

\FigureLayoutDeclare{img/2012/hunan_liberal/q4_fig1.png}{5.0cm}{102bp 0bp 240bp 66bp}
\bitmapfigure[max width=0.42\linewidth]{img/2012/hunan_liberal/q4_fig1.png}

\choices
    {\choicebitmap[width=0.15\paperwidth]{img/2012/hunan_liberal/q4_optA.png}}
    {\choicebitmap[width=0.15\paperwidth]{img/2012/hunan_liberal/q4_optB.png}}
    {\choicebitmap[width=0.15\paperwidth]{img/2012/hunan_liberal/q4_optC.png}}
    {\choicebitmap[width=0.15\paperwidth]{img/2012/hunan_liberal/q4_optD.png}}
\end{problem}

\begin{answer}
C
\end{answer}

\begin{solution}
正视图和侧视图相同，且上部矩形内部没有实线或虚线，说明上部几何体的侧棱方向与投影方向一致，斜边不能在正视图或侧视图的上部投影中出现. 若俯视图为 C，内边界含有一条斜边，该斜边对应的侧棱在正视图或侧视图中必留下斜线，与题给上部矩形内部无斜线矛盾. 因此 C 不可能，故选 C.
\end{solution}

\begin{problem}
设某大学的女生体重 \( y \) （单位： \( \mathrm{kg} \)） 与身高 \( x \) （单位： \( \mathrm{cm} \)） 具有线性相关关系，根据一组样本数据 \((x_i, y_i)\)（\(i = 1\)，\(2\)，\(\cdots\)，\(n\)），用最小二乘法建立的回归方程为 \( \hat{y} = 0.85x - 85.71 \)，则下列结论中不正确的是（\quad）

\choices
    {\(y\) 与 \(x\) 具有正的线性相关关系}
    {回归直线过样本点的中心（\(\overline{x}\)，\(\overline{y}\)）}
    {若该大学某女生身高增加 \(1 \mathrm{~cm}\)，则其体重约增加 \(0.85 \mathrm{~kg}\)}
    {若该大学某女生身高为 \(170 \mathrm{~cm}\)，则可断定其体重必为 \(58.79 \mathrm{~kg}\)}
\end{problem}

\begin{answer}
D
\end{answer}

\begin{solution}
回归方程的斜率 \(0.85>0\)，说明 \(y\) 与 \(x\) 具有正的线性相关关系，A 正确. 最小二乘回归直线必经过样本中心
\((\overline{x},\overline{y})\)，B 正确. 斜率的意义是身高每增加 \(1\mathrm{~cm}\)，按回归模型估计体重平均增加约 \(0.85\mathrm{~kg}\)，C 正确.

当 \(x=170\) 时，回归方程只给出预测值
\(\hat y=0.85\times170-85.71=58.79\)，并不能断定实际体重必为这个值. 因此不正确的是 D.
\end{solution}

\begin{problem}
已知双曲线 \(C: \frac{x^2}{a^2} - \frac{y^2}{b^2} = 1\) 的焦距为 10，点 \(P(2,1)\) 在 \(C\) 的渐近线上，则 \(C\) 的方程为（\quad）

\choices
    {\(\frac{x^2}{20} -\frac{y^2}{5} = 1\)}
    {\(\frac{x^2}{5} -\frac{y^2}{20} = 1\)}
    {\(\frac{x^2}{80} -\frac{y^2}{20} = 1\)}
    {\(\frac{x^2}{20} -\frac{y^2}{80} = 1\)}
\end{problem}

\begin{answer}
A
\end{answer}

\begin{solution}
双曲线的焦距为 \(2c=10\)，所以 \(c=5\)，从而 \(c^2=a^2+b^2=25\). 其渐近线为
\(y=\pm\frac{b}{a}x\). 点 \(P(2,1)\) 在渐近线上，故
\(\frac{b}{a}=\frac12\)，即 \(b^2=\frac{a^2}{4}\). 代入焦点关系得
\[
a^2+\frac{a^2}{4}=25
\]
所以 \(a^2=20\)，\(b^2=5\). 因而双曲线方程为
\(\frac{x^2}{20}-\frac{y^2}{5}=1\)，选 A.
\end{solution}

\begin{problem}
设 \(a > b > 1\)，\(c < 0\)，给出下列三个结论：
\begin{circlelist}
\item \( \frac{c}{a} > \frac{c}{b} \)；
\item \( a^{c} < b^{c} \)；
\item \( \log_{b}(a - c) > \log_{a}(b - c) \).
\end{circlelist}
其中所有正确结论的序号是（\quad）

\choices
    {\circled{1}}
    {\circled{1}\circled{2}}
    {\circled{2}\circled{3}}
    {\circled{1}\circled{2}\circled{3}}
\end{problem}

\begin{answer}
D
\end{answer}

\begin{solution}
因为 \(a>b>1\)，且 \(c<0\)，所以
\[
\frac{c}{a}-\frac{c}{b}=\frac{c(b-a)}{ab}>0
\]
结论一正确. 对底数大于 \(1\) 的幂函数，底数越大，负指数幂越小，因此
\(a^c<b^c\)，结论二正确.

又因为 \(a-c>b-c>1\)，并且 \(\ln a>\ln b>0\)、
\(\ln(a-c)>\ln(b-c)>0\)，于是
\[
\log_b(a-c)>\log_a(b-c)
\quad\Longleftrightarrow\quad
\frac{\ln(a-c)}{\ln b}>\frac{\ln(b-c)}{\ln a}
\]
交叉相乘后，左侧分子乘积为
\(\ln(a-c)\ln a\)，右侧为 \(\ln(b-c)\ln b\)，前者的两个正因子分别都大于后者对应因子，故结论三也正确. 因此选 D.
\end{solution}

\begin{problem}
在 \(\triangle ABC\) 中，\(AC = \sqrt{7}\)，\(BC = 2\)，\(B = 60^{\circ}\)，则 \(BC\) 边上的高等于（\quad）

\choices
    {\(\frac{\sqrt{3}}{2}\)}
    {\(\frac{3\sqrt{3}}{2}\)}
    {\(\frac{\sqrt{3} + \sqrt{6}}{2}\)}
    {\(\frac{\sqrt{3} + \sqrt{39}}{4}\)}
\end{problem}

\begin{answer}
B
\end{answer}

\begin{solution}
设 \(AB=t>0\)，以 \(B\) 为原点、\(BC\) 所在直线为 \(x\) 轴. 由
\(\angle B=60^\circ\)，点 \(A\) 的坐标可写为
\(\left(\frac t2,\frac{\sqrt3}{2}t\right)\)，而 \(C=(2,0)\). 由
\(AC=\sqrt7\) 得
\[
\left(\frac t2-2\right)^2+\left(\frac{\sqrt3}{2}t\right)^2=7
\]
即 \(t^2-2t-3=0\). 因为 \(t>0\)，所以 \(t=3\). 三角形对 \(BC\) 边的高为
\[
t\sin60^\circ=3\times\frac{\sqrt3}{2}=\frac{3\sqrt3}{2}
\]
故选 B.
\end{solution}

\begin{problem}
设定义在 \(\R\) 上的函数 \(f(x)\) 是最小正周期为 \(2\pi\) 的偶函数，\(f'(x)\) 是 \(f(x)\) 的导函数. 当 \(x \in [0, \pi]\) 时，\(0 < f(x) < 1\)；当 \(x \in (0, \pi)\) 且 \(x \neq \frac{\pi}{2}\) 时，\(\left(x - \frac{\pi}{2}\right)f'(x) > 0\)，则函数 \(y = f(x) - \sin x\) 在 \([-2\pi, 2\pi]\) 上的零点个数为（\quad）

\choices
    {2}
    {4}
    {5}
    {8}
\end{problem}

\begin{answer}
B
\end{answer}

\begin{solution}
令 \(h(x)=f(x)-\sin x\). 由题设，在 \(0<x<\frac\pi2\) 时
\(x-\frac\pi2<0\)，故 \(f'(x)<0\)，而 \(\sin x\) 递增，所以 \(h(x)\) 严格递减. 又
\(h(0)=f(0)>0\)，\(h(\frac\pi2)=f(\frac\pi2)-1<0\)，故该区间恰有一个零点.

在 \(\frac\pi2<x<\pi\) 时 \(f'(x)>0\)，而 \(\sin x\) 递减，所以 \(h(x)\) 严格递增. 由于
\(h(\frac\pi2)<0\)，\(h(\pi)=f(\pi)>0\)，该区间也恰有一个零点. 因此 \((0,\pi)\) 内有两个零点.

对 \(0<x<\pi\)，由 \(0<f(x)<1\) 和 \(\sin x>0\) 可知
\(h(-x)=f(x)+\sin x>0\)，所以 \((-\pi,0)\) 内无零点. 若 \(x\in(\pi,2\pi)\)，令 \(t=2\pi-x\in(0,\pi)\)，则
\(f(x)=f(t)\)、\(\sin x=-\sin t\)，从而 \(h(x)=f(t)+\sin t>0\). 由周期性，若 \(r\in(0,\pi)\) 是前面得到的零点，则 \(h(r-2\pi)=h(r)=0\)，所以在 \((-2\pi,-\pi)\) 内恰有对应的两个零点；正半轴的 \((\pi,2\pi)\) 及端点均无新增零点. 故 \([-2\pi,2\pi]\) 上共有 \(2+2=4\) 个零点，选 B.
\end{solution}

\section{填空题}

\begin{problem}
在极坐标系中，曲线 \(C: \rho (\sqrt{2} \cos \theta + \sin \theta) = 1\) 与曲线 \(C_{2}: \rho = a\) （\(a > 0\)） 的一个交点在极轴上，则 \(a =\fillinblank{}\).
\end{problem}

\begin{answer}
\(\dfrac{\sqrt{2}}{2}\)
\end{answer}

\begin{solution}
极轴对应 \(\theta=0\)，此时曲线 \(C\) 上的点满足
\(\rho\sqrt2=1\)，所以 \(\rho=\frac{\sqrt2}{2}\). 该点同时在
\(C_2:\rho=a\) 上，故 \(a=\frac{\sqrt2}{2}\).
\end{solution}

\begin{problem}
某制药企业为了对某种药用液体进行生物测定，需要优选培养温度，试验范围为 \(29^{\circ} \mathrm{C} \sim 63^{\circ} \mathrm{C}\)，精确度要求 \(\pm 1^{\circ} \mathrm{C}\). 用分数法进行优选时，能保证找到最佳培养温度需要的最少试验次数为\fillinblank{}.
\end{problem}

\begin{answer}
7
\end{answer}

\begin{solution}
分数法把试验区间按 Fibonacci 数列确定的试验点分割. 设
\(F_1=F_2=1\)，且 \(F_{n+1}=F_n+F_{n-1}\). 若还可进行 \(r\) 次试验，记该方法能保证分辨的等效区间份数为 \(N_r\). 一次比较后，保留的两个可能区间分别还需要 \(r-1\) 次和 \(r-2\) 次试验，因此
\(N_r=N_{r-1}+N_{r-2}\)，且 \(N_1=1\)，\(N_2=2\)，从而 \(N_r=F_{r+1}\). 所以最后包含最佳温度的区间长度至多为
\[
\frac{63-29}{F_{r+1}}=\frac{34}{F_{r+1}}
\]
精确度要求为 \(\pm1^\circ\mathrm C\)，故最后的存优区间长度应不超过 \(2^\circ\mathrm C\)，即需有
\[
F_{r+1}\geqslant\frac{34}{2}=17
\]
由于 \(F_7=13<17\)，而 \(F_8=21>17\)，所以 \(r=6\) 次不能保证达到精度，\(r=7\) 次可以保证达到精度. 因而最少试验次数为 \(7\).
\end{solution}

\begin{problem}
不等式 \(x^{2} - 5x + 6 \leqslant 0\) 的解集为\fillinblank{}.
\end{problem}

\begin{answer}
\([2,3]\)
\end{answer}

\begin{solution}
因式分解得
\[
x^2-5x+6=(x-2)(x-3)
\]
抛物线开口向上，乘积不大于零正好发生在两根之间，包含两个根，因此解集为
\([2,3]\).
\end{solution}

\begin{problem}
如图是某学校一名篮球运动员在五场比赛中所得分数的茎叶图，则该运动员在这五场比赛中得分的方差为\fillinblank{}.
\begin{center}
\begin{tabular}{c|ccc}
\(0\) & \(8\) & \(9\) & \\
\(1\) & \(0\) & \(3\) & \(5\)
\end{tabular}
\end{center}
（注：方差 \( s^{2}=\frac{1}{n}\left[(x_{1}-\overline{x})^{2}+(x_{2}-\overline{x})^{2}+\cdots+(x_{n}-\overline{x})^{2}\right] \)，其中 \( \overline{x} \) 为 \(x_{1}\)，\(x_{2}\)，\(\cdots\)，\(x_{n}\) 的平均数）
\end{problem}

\begin{answer}
6.8
\end{answer}

\begin{solution}
由茎叶图可知五场得分为 \(8\)，\(9\)，\(10\)，\(13\)，\(15\). 平均数为
\[
\overline{x}=\frac{8+9+10+13+15}{5}=11
\]
因此
\[
s^2=\frac{(8-11)^2+(9-11)^2+(10-11)^2+(13-11)^2+(15-11)^2}{5}
 =\frac{9+4+1+4+16}{5}=6.8
\]
\end{solution}

\begin{problem}
如果执行如图所示的程序框图，输入 \(x = 4.5\)，则输出的数 \(i =\fillinblank{}\).
\texfigure[max width=0.65\linewidth]{img_repaint/2012/hunan_liberal/q14_fig1.tex}
\end{problem}

\begin{answer}
4
\end{answer}

\begin{solution}
程序先令 \(i=1\)，每次循环先执行 \(x=x-1\)，再判断是否有 \(x<1\). 输入 \(x=4.5\) 时各次判断如下：
\(4.5\to3.5\)，\((i=1)\)，\(3.5\to2.5\)，\((i=2)\)，\(2.5\to1.5\)，\((i=3)\)
前三次判断均为否，每次再令 \(i=i+1\). 第四次得到
\(1.5\to0.5\)，此时 \(x<1\)，不再增加 \(i\)，输出 \(i=4\).
\end{solution}

\begin{problem}
如图，在平行四边形 \(ABCD\) 中，\(AP \perp BD\)，垂足为 \(P\)，且 \(AP = 3\)，则 \(\overrightarrow{AP} \cdot \overrightarrow{AC} =\)\fillinblank{}.
\bitmapfigure[max width=0.62\linewidth]{img/2012/hunan_liberal/q15_fig1.png}
\end{problem}

\begin{answer}
18
\end{answer}

\begin{solution}
取 \(P\) 为原点，令 \(BD\) 为 \(x\) 轴. 由于 \(AP\perp BD\) 且 \(AP=3\)，可设
\(A=(0,3)\)、\(B=(-u,0)\)、\(D=(v,0)\)，其中 \(u\)，\(v>0\). 平行四边形的对角线互相平分，故
\(A+C=B+D\)，从而 \(C=(v-u,-3)\). 因此
\(\overrightarrow{AP}=(0,-3)\)，\(\overrightarrow{AC}=(v-u,-6)\)
于是
\[
\overrightarrow{AP}\cdot\overrightarrow{AC}=0\cdot(v-u)+(-3)(-6)=18
\]
\end{solution}

\begin{problem}
对于 \(n \in \mathbf{N}^*\)，将 \(n\) 表示为 \(n = a_k \times 2^k + a_{k-1} \times 2^{k-1} + \cdots + a_1 \times 2^1 + a_0 \times 2^0\)，当 \(i = k\) 时，\(a_i = 1\)，当 \(0 \leqslant i \leqslant k - 1\) 时，\(a_i\) 为 0 或 1. 定义 \(b_n\) 如下：在 \(n\) 的上述表示中，当 \(a_0\)，\(a_1\)，\(a_2\)，\(\cdots\)，\(a_k\) 中等于 1 的个数为奇数时，\(b_n = 1\)；否则 \(b_n = 0\).
\begin{enumerate}
\item \( b_{2}+b_{4}+b_{6}+b_{8}=\fillinblank{}\)；
\item 记 \(c_{m}\) 为数列 \(\{b_{n}\}\) 中第 \(m\) 个为 0 的项与第 \(m + 1\) 个为 0 的项之间的项数，则 \(c_{m}\) 的最大值是\fillinblank{}.
\end{enumerate}
\end{problem}

\begin{answer}
\(3\)；\(2\)
\end{answer}

\begin{solution}
把 \(n\) 写成二进制后，\(b_n\) 等于二进制表示中 \(1\) 的个数的奇偶性指标. 于是
\(2=10_2\)、\(4=100_2\)、\(6=110_2\)、\(8=1000_2\)，分别有
\(b_2=1\)，\(b_4=1\)，\(b_6=0\)，\(b_8=1\)，故
\[
b_2+b_4+b_6+b_8=1+1+0+1=3
\]

设 \(s(n)\) 为 \(n\) 的二进制表示中 \(1\) 的个数. 由二进制移位可得
\(s(2k)=s(k)\)，\(s(2k+1)=s(k)+1\).
对 \(k\geqslant1\)，有 \(b_{2k}=b_k\)，\(b_{2k+1}=1-b_k\). 若三个连续项从偶数项
\(2k\) 开始，则 \(b_{2k}\) 与 \(b_{2k+1}\) 一定一为 \(0\)；若从奇数项 \(2k+1\) 开始且
\(k\geqslant1\)，则 \(b_{2k+2}\) 与 \(b_{2k+3}\) 一定一为 \(0\). 起始的三个项为
\(b_1=1\)，\(b_2=1\)，\(b_3=0\)，也不可能连续出现三个 \(1\). 所以连续的 \(1\) 最多两个.

又 \(b_6=0\)，\(b_7=1\)，\(b_8=1\)，\(b_9=0\)，所以两个相邻的零项之间可以有两项，且不能更多，故 \(c_m\) 的最大值为 \(2\).
\end{solution}

\section{解答题}

\begin{problem}
某超市为了解顾客的购物量及结算时间等信息，安排一名员工随机收集了在该超市购物的 100 位顾客的相关数据，如下表所示.

\begin{center}
\begin{tabular}{c|ccccc}
一次购物量 & \(1\text{\ 至\ }4\)件 & \(5\text{\ 至\ }8\)件 & \(9\text{\ 至\ }12\)件 & \(13\text{\ 至\ }16\)件 & \(17\)件及以上 \\
\hline
顾客数（人） & \(x\) & 30 & 25 & \(y\) & 10 \\
\hline
结算时间（分钟/人） & 1 & 1.5 & 2 & 2.5 & 3
\end{tabular}
\end{center}

已知这 100 位顾客中一次购物量超过 8 件的顾客占 55\%.
\begin{enumerate}
\item 确定 \(x\)，\(y\) 的值，并估计顾客一次购物的结算时间的平均值；
\item 求一位顾客一次购物的结算时间不超过 2 分钟的概率. （将频率视为概率）
\end{enumerate}

\end{problem}

\begin{solution}
\begin{enumerate}
\item 购物量超过 \(8\) 件的顾客数为 \(25+y+10\)，由题意
\[
25+y+10=100\times55\%=55
\]
所以 \(y=20\). 再由总顾客数为 \(100\)，得
\[
x+30+25+y+10=100
\]
\(y=20\) 代入上式，得 \(x=15\).

将五组顾客数分别乘以对应的结算时间再求平均，得到平均结算时间的估计值
\[
\frac{15\times1+30\times1.5+25\times2+20\times2.5+10\times3}{100}
=\frac{190}{100}=1.9\text{ 分钟}
\]
\item 结算时间不超过 \(2\) 分钟的顾客来自前三组，因此所求概率为
\[
\frac{15+30+25}{100}=\frac{70}{100}=\frac{7}{10}
\]
\end{enumerate}
\end{solution}

\begin{problem}
已知函数 \( f(x) = A \sin (\omega x + \varphi) \) （\(x \in \R\)，\(\omega > 0\)，\(0 < \varphi < \tfrac{\pi}{2}\)）的部分图象如图所示.

\begin{enumerate}
\item 求函数 \( f(x) \) 的解析式；
\item 求函数 \( g(x)=f\left(x-\frac{\pi}{12}\right)-f\left(x+\frac{\pi}{12}\right) \) 的单调递增区间.
\end{enumerate}
\bitmapfigure[max width=0.58\linewidth]{img/2012/hunan_liberal/q18_fig1.png}

\end{problem}

\begin{solution}
\begin{enumerate}
\item 由图象可知 \(f(0)=1\)，且从 \(x=0\) 向右的第一个零点为
\(x=\frac{5\pi}{12}\). 因为 \(0<\varphi<\frac{\pi}{2}\)，该零点对应正弦函数自相位
\(\varphi\) 后的第一次过零，即
\[
2\cdot\frac{5\pi}{12}+\varphi=\pi
\]
所以 \(\varphi=\frac{\pi}{6}\). 图象上相邻两个零点为
\(\frac{5\pi}{12}\) 和 \(\frac{11\pi}{12}\)，相距 \(\frac{\pi}{2}\)，这是半个周期，故
\(T=\pi\)，\(\omega=\frac{2\pi}{T}=2\).
再由 \(f(0)=A\sin\varphi=1\)，得到
\[
A\sin\frac{\pi}{6}=1
\]
从而 \(A=2\). 因此
\[
f(x)=2\sin\left(2x+\frac{\pi}{6}\right)
\]
\item 代入 \(f\) 化简：
\[
g(x)=2\sin\left(2x\right)-2\sin\left(2x+\frac{\pi}{3}\right)
=-2\cos\left(2x+\frac{\pi}{6}\right)
\]
所以
\[
g'(x)=4\sin\left(2x+\frac{\pi}{6}\right)
\]
函数 \(g\) 在区间上单调递增需要 \(g'(x)\geqslant0\)，即
\(2k\pi\leqslant2x+\frac{\pi}{6}\leqslant(2k+1)\pi\)，\(k\in\Z\).
解得单调递增区间为
\(\left[k\pi-\frac{\pi}{12},\,k\pi+\frac{5\pi}{12}\right]\)，\(k\in\Z\).
\end{enumerate}
\end{solution}

\begin{problem}
如图，在四棱锥 \(P - ABCD\) 中，\(PA \perp\) 平面 \(ABCD\)，底面 \(ABCD\) 是等腰梯形，\(AD\parallel BC\)，\(AC \perp BD\).

\begin{enumerate}
\item 证明： \(BD \perp PC\)；
\item 若 \(AD=4\)，\(BC=2\)，直线 \(PD\) 与平面 \(PAC\) 所成的角为 \(30^{\circ}\)，求四棱锥 \(P-ABCD\) 的体积.
\end{enumerate}
\bitmapfigure[max width=0.62\linewidth]{img/2012/hunan_liberal/q19_fig1.png}

\end{problem}

\begin{solution}
\begin{enumerate}
\item 因为 \(PA\perp\) 平面 \(ABCD\)，而 \(BD\) 在平面 \(ABCD\) 内，所以
\(BD\perp PA\). 又已知 \(AC\perp BD\). 直线 \(PA\)、\(AC\) 是平面 \(PAC\) 内相交的两条直线，
因此 \(BD\perp\) 平面 \(PAC\)，从而 \(BD\perp PC\).
\item 设 \(PA=h\). 取底面所在平面为 \(z=0\)，令等腰梯形的两条底边水平，设
\(A=(-2,0,0)\)，\(D=(2,0,0)\)，\(B=(-1,t,0)\)，\(C=(1,t,0)\).
这里使用了等腰梯形的对称性，且 \(AD=4\)，\(BC=2\). 由 \(AC\perp BD\)，有
\(\overrightarrow{AC}=(3,t,0)\)，\(\overrightarrow{BD}=(3,-t,0)\)
故 \(9-t^2=0\). 因为 \(t>0\)，所以 \(t=3\). 底面积为
\[
S_{ABCD}=\frac{AD+BC}{2}\,t=\frac{4+2}{2}\times3=9
\]
令 \(P=(-2,0,h)\). 由第（1）问，\(BD\) 是平面 \(PAC\) 的一个法向方向. 于是直线 \(PD\) 与平面 \(PAC\) 所成角 \(\theta\) 满足
\[
\sin\theta
=\frac{\left|\overrightarrow{DP}\cdot\overrightarrow{DB}\right|}
{|\overrightarrow{DP}|\,|\overrightarrow{DB}|}
\]
其中
\(\overrightarrow{DP}=(-4,0,h)\)，\(\overrightarrow{DB}=(-3,3,0)\)
所以
\[
\sin30^\circ
=\frac{12}{\sqrt{16+h^2}\cdot3\sqrt2}
\]
解得
\(\frac12=\frac{4}{\sqrt{2}\sqrt{16+h^2}}\)，\(h^2=16\).
由于 \(h>0\)，得 \(h=4\). 因而四棱锥体积为
\[
V=\frac13S_{ABCD}\,PA=\frac13\times9\times4=12
\]
\end{enumerate}
\end{solution}

\begin{problem}
某公司一下属企业从事某种高科技产品的生产. 该企业第一年年初有资金2000万元，将其投入生产，到当年年底资金增长了 \(50\%\)，预计以后每年资金年增长率与第一年的相同. 公司要求企业从第一年开始，每年年底上缴资金 \(d\) 万元，并将剩余资金全部投入下一年生产. 设第 \(n\) 年年底企业上缴资金后的剩余资金为 \(a_{n}\) 万元.
\begin{enumerate}
\item 用 \(d\) 表示 \(a_{1}\)，\(a_{2}\)，并写出 \( \{a_{n+1}\} \) 与 \( a_{n} \) 的关系式；
\item 若公司希望经过 \(m\) 年（\(m \geqslant 3\)）使企业的剩余资金为 \(4000\) 万元，试确定企业每年上缴资金 \(d\) 的值（用 \(m\) 表示）.
\end{enumerate}

\end{problem}

\begin{solution}
\begin{enumerate}
\item 第一年年初的资金为 \(2000\) 万元，年底增长 \(50\%\) 后为
\(2000\times1.5=3000\) 万元，缴纳 \(d\) 万元后
\[
a_1=3000-d
\]
第二年年底上缴前资金为 \(1.5a_1\)，所以
\[
a_2=\frac32a_1-d
\]
第 \(n\) 年年底剩余资金为 \(a_n\) 万元，下一年增长 \(50\%\) 后为 \(\frac32a_n\) 万元，再上缴 \(d\) 万元，所以对 \(n\geqslant1\)，有
\[
a_{n+1}=\frac32a_n-d
\]
\item 递推式的固定值为 \(2d\)，因为
\(2d=\frac32(2d)-d\). 因此
\[
a_{n+1}-2d=\frac32(a_n-2d)
\]
从而
\[
a_m-2d=\left(\frac32\right)^{m-1}(a_1-2d)
=\left(\frac32\right)^{m-1}(3000-3d)
\]
由 \(a_m=4000\)，得
\[
4000=2d+3000\left(\frac32\right)^{m-1}
-3d\left(\frac32\right)^{m-1}
\]
整理并乘以 \(2^{m-1}\)，得到
\[
d=\frac{3000\cdot3^{m-1}-4000\cdot2^{m-1}}{3^m-2^m}
=\frac{1000\left(3^m-2^{m+1}\right)}{3^m-2^m}
\]
\end{enumerate}
\end{solution}

\begin{problem}
在直角坐标系 \(xOy\) 中，已知中心在原点，离心率为 \(\frac{1}{2}\) 的椭圆 \(E\) 的一个焦点为圆 \(C: x^2 + y^2 - 4x + 2 = 0\) 的圆心.
\begin{enumerate}
\item 求椭圆 \(E\) 的方程；
\item 设 \(P\) 是椭圆 \(E\) 上一点，过 \(P\) 作两条斜率之积为 \(\frac{1}{2}\) 的直线 \(l_{1}\)，\(l_{2}\)，当直线 \(l_{1}\)，\(l_{2}\) 都与圆 \(C\) 相切时，求 \(P\) 的坐标.
\end{enumerate}

\end{problem}

\begin{solution}
\begin{enumerate}
\item 将圆 \(C\) 的方程配方：
\[
(x-2)^2+y^2=2
\]
所以圆心为 \((2,0)\). 该圆心是椭圆的一个焦点，故椭圆焦距参数
\(c=2\). 由离心率 \(e=\frac ca=\frac12\)，得 \(a=4\). 再由
\(a^2=b^2+c^2\)，得 \(b^2=16-4=12\). 因此椭圆方程为
\[
\frac{x^2}{16}+\frac{y^2}{12}=1
\]
\item 设 \(P=(x,y)\)，过 \(P\) 的斜率为 \(k\) 的直线为
\[
kX-Y+y-kx=0
\]
它与圆 \(C\) 相切的条件是圆心到直线的距离等于半径 \(\sqrt2\)，即
\[
\frac{\left|2k+y-kx\right|}{\sqrt{k^2+1}}=\sqrt2
\]
平方并整理，得到关于切线斜率 \(k\) 的方程
\[
\bigl((2-x)^2-2\bigr)k^2+2y(2-x)k+y^2-2=0
\]
题设两条切线的斜率之积为 \(\frac12\)，由韦达定理
\[
\frac{y^2-2}{(2-x)^2-2}=\frac12
\]
所以
\[
x^2-4x-2y^2+6=0
\]
另一方面 \(P\) 在椭圆上，故
\[
3x^2+4y^2=48
\]
将前一式乘以 \(2\)，并消去 \(4y^2\)，得
\[
5x^2-8x-36=0
\]
于是
\[
x=-2\quad\text{或}\quad x=\frac{18}{5}
\]
当 \(x=-2\) 时，由椭圆方程得 \(y^2=9\)，即 \(y=\pm3\). 当
\(x=\frac{18}{5}\) 时，得 \(y^2=\frac{57}{25}\)，即
\(y=\pm\frac{\sqrt{57}}5\). 所以
\(P=(-2,3)\)，\(\ (-2,-3)\)，\(\ \left(\frac{18}{5},\frac{\sqrt{57}}5\right)\)，\(\ \left(\frac{18}{5},-\frac{\sqrt{57}}5\right)\).
这些点到圆心的距离平方分别为 \(25\) 和 \(\frac{121}{25}\)，均大于圆半径平方 \(2\)，确实各有两条实切线；且 \((2-x)^2-2\neq0\)，斜率均为有限值，满足题设条件.
\end{enumerate}
\end{solution}

\begin{problem}
已知函数 \( f(x) = \e^{x} - ax \)，其中 \( a > 0 \).
\begin{enumerate}
\item 若对一切 \(x \in \R\)，\(f(x) \geqslant 1\) 恒成立，求 \(a\) 的取值集合；
\item 在函数 \(f(x)\) 的图象上取定两点 \(A(x_{1}, f(x_{1}))\)，\(B(x_{2}, f(x_{2}))\)（\(x_{1} < x_{2}\)），记直线 \(AB\) 的斜率为 \(k\)，证明：存在 \(x_{0} \in (x_{1}, x_{2})\)，使 \(f'(x_{0}) = k\) 成立.
\end{enumerate}

\end{problem}

\begin{solution}
\begin{enumerate}
\item 函数
\(f'(x)=\e^x-a\)，\(f''(x)=\e^x>0\).
所以 \(f\) 在 \(\R\) 上有唯一极小值点
\(x=\ln a\)，且该点的函数值为
\[
f(\ln a)=\e^{\ln a}-a\ln a=a(1-\ln a)
\]
令 \(t=\ln a\)，则极小值为 \(e^t(1-t)\). 对其求导得
\[
\frac{\mathrm d}{\mathrm dt}\bigl(e^t(1-t)\bigr)=-te^t
\]
因此它在 \(t<0\) 时递增，在 \(t>0\) 时递减，最大值为
\(e^0(1-0)=1\)，且只有 \(t=0\) 时取到. 题设要求对一切
\(x\) 有 \(f(x)\geqslant1\)，故极小值必须等于 \(1\)，从而
\(t=0\)，即 \(a=1\). 因此 \(a\) 的取值集合为 \(\{1\}\).
\item 函数 \(f(x)=\e^x-ax\) 在闭区间 \([x_1,x_2]\) 上连续，在开区间
\((x_1,x_2)\) 上可导. 直线 \(AB\) 的斜率为
\[
k=\frac{f(x_2)-f(x_1)}{x_2-x_1}
\]
由拉格朗日中值定理，存在 \(x_0\in(x_1,x_2)\)，使得
\[
f'(x_0)=\frac{f(x_2)-f(x_1)}{x_2-x_1}=k
\]
这就证明了所要求的结论.
\end{enumerate}
\end{solution}
